\documentclass[a4paper]{article}
\usepackage[pdftex]{graphicx}
\usepackage[utf8]{inputenc}
\usepackage{enumerate}
\usepackage{siunitx}
\sisetup{locale=DE}
\usepackage{href-ul}
\hypersetup{
	colorlinks=true,
	linkcolor=blue,
	urlcolor=blue}
\usepackage{icomma}
\usepackage{amssymb}
\usepackage{geometry}
\geometry{a4paper, top=15mm, left=15mm, right=15mm, bottom=15mm,
	headsep=10mm, footskip=12mm}

\begin{document}
	%Title
	{\bf Tasks conductor characteristics }
	
	\begin{enumerate}[1.]
		
		%Task 1
		\begin{minipage}{0.65\textwidth}
			{\bf \item In an experiment, the current intensity $I$ is examined as a function of the voltage $U$ }for a constantan wire, a graphite pencil lead and an iron wire. The following diagram results.
		\end{minipage}
		\hfill
		\begin{minipage}{0.3\textwidth}
			\includegraphics[width=4 cm]{leiken185}
		\end{minipage}
		
		\begin{enumerate}[a)]
			\item What can be said about the resistance values of the three conductors as the voltage increases? Match the three graphs with the corresponding conductors. What is the general name for such graphs?
			\item Which of the three materials conducts best at a voltage of $\SI{4.0}{\volt}$?
			Calculate the associated resistance value for this conductor.
			\item Use the particle model to explain the course of the characteristic curve of conductor 2.
			\item Explain the terms PTC and NTC and what a cold conductor and what a hot conductor is.
		\end{enumerate}
		
		%Exercise 2
		{\bf \item In an experiment, the current intensity $I$ is examined as a function of the voltage $U$ for two different electrical conductors.} The following measured values result:
		
		\begin{tabular}{cccccccccc}
			& $U$ in $\SI{}{\volt}$ & 0 & 1.8 & 3.6 & 5.4 & 7.2 & 9.0 & 10.8 \\
			\hline
			Conductor 1 & $I$ in $\SI{}{\ampere}$ & 0 & 0.31 & 0.55 & 0.74 & 0.84 & 0.93 & 0.96 \\
			Conductor 2 & $I$ in $\SI{}{\ampere}$ & 0 & 0.18 & 0.37 & 0.53 & 0.71 & 0.92 & 1.06
		\end{tabular}
		
		\begin{enumerate}[a)]
			\item Evaluate the measurement table for conductor 1 numerically and formulate the test result.
			\item For conductor 2, graphically represent the current intensity as a function of the voltage in a diagram, formulate and justify the test result.
			\item What statements can be made about the resistance values of the two conductors with the help of the numerical and graphical evaluation?
		\end{enumerate}
		
		%Task 3
		\item {\bf Electrical resistance }
		
		\begin{enumerate}[a)]
			\item The electrical resistance of a light bulb is to be determined experimentally. To do this, an applied voltage of $\SI{6.0}{\volt}$ and a current of $\SI{5.0}{\ampere}$ are measured. Calculate the resistance of the conductor.
			\item The resistance of a conductor is $\SI{50}{\ohm}$. The current in the wire is determined to be $\SI{380}{\milli\ampere}$. What voltage is present?
			\item A light bulb with power $P=\SI{40}{\watt}$ is connected to the household network ($U=\SI{230}{\volt}$). Calculate the operating current and the operating resistance.
			\item A current intensity of $\SI{15}{\milli\ampere}$ can be fatal to humans. Under unfavorable conditions, the electrical resistance of the human body is approximately
			$\SI{1,0}{\kilo\ohm}$. What voltage can be life-threatening under these circumstances?
		\end{enumerate}
	\end{enumerate}
	
	\fbox{
		\begin{minipage}{0.5\textwidth}
			For the solution please \href{https://www.okuyakl.de/phys/p8widL185/le185.pdf}{click here} or scan the QR code.\\
			Additional worksheets are available at
			
			\href{http://www.okuyakl.com}{www.okuyakl.com}
		\end{minipage}
		\hfill
		\begin{minipage}{0.4\textwidth}
			\includegraphics[width=1.5 cm]{../../viecher/zwe03}
			\includegraphics[width=3 cm]{qre185}
			\includegraphics[width=2 cm]{../../viecher/afanticon1}
			
	\end{minipage}}
	
\end{document}